\chapter{Internship: Pengembangan Fitur dan Modul Proyek}

\section{Identitas Mata Kuliah}
\begingroup
\renewcommand{\arraystretch}{1.15}
\noindent\begin{tabular}{@{}p{3.6cm}@{\hspace{0.2cm}}p{0.2cm}@{\hspace{0.2cm}}p{7.6cm}@{}}
    \textbf{Kode Mata Kuliah} & : & MII21-3011 \\
    \textbf{Nama Mata Kuliah} & : & Internship: Pengembangan Fitur dan Modul Proyek \\
    \textbf{Bobot} & : & 4 SKS \\
    \textbf{Bentuk Kegiatan} & : & Magang MBKM sebagai Software Engineer Intern di OPPO Manufacturing Indonesia \\
    \textbf{Proyek Terkait} & : & 360 Peer-to-Peer Grading System dan PQPR Software Project \\
\end{tabular}
\par\vspace{0.5\baselineskip}
\endgroup

Mata kuliah ini berfokus pada penerapan kemampuan pengembangan fitur dan modul perangkat lunak dalam lingkungan kerja industri. Selama kegiatan magang, capaian mata kuliah diwujudkan melalui kontribusi pada dua proyek internal perusahaan, yaitu 360 Peer-to-Peer Grading System dan PQPR Software Project. Kedua proyek tersebut menuntut pemahaman kebutuhan pengguna, perancangan alur kerja aplikasi, implementasi fitur, integrasi data, serta penyesuaian modul agar dapat digunakan dalam proses bisnis perusahaan.

\section{Silabus}
Silabus mata kuliah mencakup pengembangan fitur perangkat lunak berdasarkan kebutuhan pengguna, perancangan modul aplikasi, implementasi antarmuka dan logika bisnis, integrasi dengan sumber data, serta dokumentasi hasil pengembangan. Dalam konteks magang, silabus tersebut diterapkan melalui pengembangan sistem evaluasi kinerja antarkaryawan pada proyek 360 Peer-to-Peer Grading System dan pengembangan modul pendukung optimasi perubahan lini produksi pada PQPR Software Project.

Kegiatan pengembangan fitur tidak hanya menekankan penulisan kode, tetapi juga kemampuan menerjemahkan kebutuhan bisnis menjadi rancangan teknis yang dapat diimplementasikan. Setiap fitur perlu memiliki tujuan yang jelas, alur penggunaan yang dapat dipahami, dan keluaran yang sesuai dengan kebutuhan operasional pengguna.

\section{Aktivitas}
Pada proyek 360 Peer-to-Peer Grading System, aktivitas pengembangan diarahkan pada penyusunan fitur yang mendukung proses penilaian rekan kerja secara terstruktur. Fitur yang dikembangkan mencakup alur pengisian penilaian, pengelolaan data evaluasi, tampilan rekapitulasi, serta mekanisme yang membantu pengguna menyelesaikan proses penilaian dengan konsisten. Pengembangan dilakukan dengan memperhatikan kemudahan penggunaan karena sistem ditujukan untuk pengguna dari berbagai tingkat jabatan dan latar belakang teknis.

Pada PQPR Software Project, aktivitas pengembangan berfokus pada modul yang membantu proses changeover produk di lingkungan manufaktur. Modul tersebut berkaitan dengan pengelolaan informasi produk, penyusunan konfigurasi lini produksi, serta pemanfaatan data kemampuan operator untuk mendukung keputusan operasional. Pengembangan fitur pada proyek ini membutuhkan pemahaman terhadap alur kerja pabrik, hubungan antar data produksi, dan kebutuhan pembaruan informasi secara akurat.

Selain implementasi fitur, aktivitas yang dilakukan juga meliputi diskusi kebutuhan dengan pembimbing dan anggota tim, pembacaan spesifikasi, penyesuaian rancangan berdasarkan umpan balik, serta perbaikan modul setelah dilakukan pengecekan. Aktivitas tersebut memperlihatkan bahwa pengembangan perangkat lunak di industri berlangsung secara iteratif dan membutuhkan koordinasi lintas peran.

\section{Konsep Dasar yang Berkaitan dengan Silabus}
Konsep utama yang digunakan dalam mata kuliah ini adalah pemecahan kebutuhan bisnis menjadi modul perangkat lunak yang memiliki tanggung jawab jelas. Modul yang baik perlu memiliki batas fungsi yang terdefinisi, masukan dan keluaran yang dapat dipahami, serta perilaku yang konsisten ketika digunakan bersama modul lain. Prinsip ini penting pada proyek 360 Peer-to-Peer Grading System karena alur penilaian, pengolahan data, dan penyajian hasil perlu saling terhubung tanpa menimbulkan kebingungan bagi pengguna.

Konsep lain yang digunakan adalah perancangan berbasis kebutuhan pengguna. Fitur tidak cukup hanya dapat berjalan secara teknis, tetapi harus sesuai dengan cara pengguna menyelesaikan pekerjaannya. Pada proyek PQPR Software Project, rancangan modul perlu mempertimbangkan konteks kerja manufaktur, seperti kebutuhan informasi yang cepat, data yang akurat, serta perubahan kondisi produksi yang dapat terjadi sewaktu-waktu.

Pengembangan fitur juga berkaitan dengan integrasi data dan validasi alur. Data yang dimasukkan, diproses, dan ditampilkan harus memiliki struktur yang konsisten agar modul dapat digunakan sebagai dasar pengambilan keputusan. Oleh karena itu, setiap rancangan fitur perlu memperhatikan relasi data, pembatasan akses, validasi masukan, dan kejelasan tampilan keluaran.

\section{Hasil dan Pembahasan}
Hasil utama dari pelaksanaan mata kuliah ini adalah meningkatnya kemampuan menerapkan proses pengembangan fitur secara utuh, mulai dari memahami kebutuhan hingga menghasilkan modul yang dapat digunakan dalam sistem internal perusahaan. Pada proyek 360 Peer-to-Peer Grading System, pengembangan fitur mendukung proses evaluasi kinerja yang lebih terstruktur dan terdokumentasi. Sistem membantu proses pengumpulan umpan balik sehingga data penilaian dapat dikelola dengan lebih rapi dibandingkan proses manual.

Pada PQPR Software Project, modul yang dikembangkan mendukung kebutuhan perusahaan dalam mengelola informasi produksi dan konfigurasi changeover. Dengan adanya modul yang terintegrasi, informasi produk dan kemampuan operator dapat digunakan untuk membantu penyusunan konfigurasi lini yang lebih efisien. Hal ini menunjukkan keterkaitan langsung antara pengembangan perangkat lunak dan peningkatan efektivitas proses operasional di lingkungan manufaktur.

Pembahasan dari kegiatan ini menunjukkan bahwa pengembangan fitur di industri memerlukan keseimbangan antara aspek teknis dan pemahaman domain. Keputusan teknis perlu mempertimbangkan dampaknya terhadap pengguna, pemeliharaan sistem, dan kesesuaian dengan proses bisnis. Umpan balik dari pembimbing dan anggota tim menjadi bagian penting karena membantu memastikan bahwa fitur yang dibuat tidak hanya selesai secara implementasi, tetapi juga relevan dengan kebutuhan perusahaan.

\section{Kesimpulan}
Mata kuliah Internship: Pengembangan Fitur dan Modul Proyek terlaksana melalui kontribusi pada pengembangan 360 Peer-to-Peer Grading System dan PQPR Software Project. Kegiatan tersebut memberikan pengalaman dalam menerjemahkan kebutuhan bisnis menjadi fitur perangkat lunak, merancang modul yang terstruktur, melakukan integrasi data, serta menyesuaikan hasil pengembangan berdasarkan umpan balik. Pengalaman ini memperkuat pemahaman bahwa keberhasilan pengembangan fitur tidak hanya ditentukan oleh kemampuan teknis, tetapi juga oleh ketepatan memahami kebutuhan pengguna dan konteks operasional sistem.
